// 323111047

\section{Introducción}

Planteamiento del problema. El reto técnico de esta actividad consiste en diseñar una aplicación en Python 3 que resuelva el problema de selección de actividades aplicado a un horario académico dinámico. El problema central radica en identificar un conjunto máximo de materias que puedan cursarse sin que existan traslapes de horario, considerando que cada asignatura posee una hora de inicio, una de fin y un nivel de prioridad asignado por el usuario. Para resolver esto, se implementa una solución que combina una estrategia ávida mediante el algoritmo de ordenamiento por inserción para priorizar las materias que terminan más temprano, y una estrategia incremental que construye el horario validando conflictos de tiempo y resolviendo empates mediante criterios de prioridad o interacción directa con el usuario.

\vspace{0.4cm}

Motivación. La importancia de realizar este ejercicio se basa en la optimización de recursos y la gestión eficiente del tiempo, problemas que son fundamentales en la ingeniería de software y la investigación de operaciones. Comprender cómo los algoritmos ávidos pueden ofrecer soluciones óptimas de manera rápida es esencial para el desarrollo de sistemas de logística y asignación de tareas. Esta práctica motiva al estudiante a analizar la eficiencia de su código mediante el estudio de la complejidad en tiempo y espacio, permitiendo transitar de una programación empírica a un diseño algorítmico formal que utiliza estructuras de datos modernas como diccionarios para representar entidades complejas.

\vspace{0.4cm}

Objetivos. Se busca que el alumno demuestre dominio en la aplicación de criterios de ordenamiento y en la validación incremental de datos para asegurar la integridad de la información procesada. Finalmente, se pretende realizar un análisis de complejidad que verifique la eficiencia del uso de insertion sort y el recorrido lineal del horario, proporcionando una base sólida para discutir los resultados obtenidos y la efectividad de la resolución de conflictos por prioridad en la conclusión del reporte.